% ============================================================
% Artigo do Projeto AV2 — GXX (preencher)
% Compilar: latexmk -pdf artigo.tex
% Regra: tabelas e figuras entram por \input{} / \includegraphics{}
%        do que os scripts geraram. Nunca digite números à mão.
% ============================================================
\documentclass[11pt,a4paper]{article}

\usepackage[utf8]{inputenc}
\usepackage[T1]{fontenc}
\usepackage[brazil]{babel}
\usepackage{graphicx}
\usepackage{booktabs}
\usepackage{amsmath}
\usepackage{siunitx}
\usepackage{hyperref}
\usepackage[margin=2.5cm]{geometry}

\title{preencher: título do artigo (ODS + decisão)}
\author{GXX --- preencher integrantes \\ Machine Learning --- FACAPE 2026.2}
\date{\today}

\begin{document}
\maketitle

\begin{abstract}
% 120-180 palavras: problema, dados, método, resultado com número, decisão.
preencher.
\end{abstract}

\section{Introdução}
% ODS + meta numerada; problema local; pergunta; decisão; contribuição.
preencher.

\section{Métodos}
\subsection{Dados}
% fonte, licença, recorte, volume, alvo, limpeza, LGPD.
preencher.

\subsection{Preparação e validação}
% pipeline, split, CV repetida, ausência de vazamento.
preencher.

\subsection{Modelos e protocolo estatístico}
% baseline + 3 famílias; folds idênticos; teste escolhido; efeito e IC.
preencher.

\section{Resultados}
% \input{../reports/tables/tabela-comparativa.tex}
% \includegraphics[width=\linewidth]{../reports/figures/figura-comparacao.png}
preencher.

\section{Discussão}
% o que o modelo erra, por subgrupo, interpretação, ameaças à validade.
preencher.

\section{Conclusão}
% resposta à pergunta, decisão recomendada, limitações, trabalho futuro.
preencher.

\section*{Disponibilidade}
Código, dados (link/hash) e instruções: preencher (URL do repositório, tag
\texttt{v1.0-entrega}, comando \texttt{docker run}).

\bibliographystyle{plain}
\bibliography{referencias}

\end{document}
